% ---------------------------------------------------------------------
%  Chapter 15 : The Spatial Heat-Kernel Propagator and the Analytic IFT
% ---------------------------------------------------------------------
\chapter{The Spatial Heat-Kernel Propagator and the Analytic IFT}\label{ch:spatial-heatkernel}

Every chapter before this one has worked, ultimately, in a \emph{mixed}
representation. The temporal pipeline lives in frequency $\omega$; the spatial
core that precedes this chapter keeps a formal Laplacian operator symbol
floating through the algebra, untouched, because the Laplacian cannot be folded
into the $\omega$-pole finder the way a mass or a damping rate can. At some
point that mixed representation has to be cashed out into a physical answer ---
a correlation function $C(x,\tau)$ you can plot against a separation $x$ and a
time lag $\tau$. This chapter is the layer that does the cashing out.

The whole subsystem rests on a single, very old idea: \emph{the inverse Fourier
transform of a Gaussian is a Gaussian}. For a diffusive (reaction--diffusion,
Allen--Cahn-like) field the momentum-space propagator is, after one contour
integral, a Gaussian in the wavevector $k$. Its inverse transform back to real
space is therefore another Gaussian --- the \term{heat kernel}, the Green's
function of the diffusion equation $\partial_t \phi = B\,\Lap\phi$. The payoff
of recognising this is enormous: the spatial transform is done \emph{with
pencil and paper}, in closed form, and never as a discrete numerical transform
on a grid of momenta. The slogan inside the codebase is that the analytic heat
kernel ``retired the $q$-grid.''

We will build this from the ground up. First the physics: the inverse
propagator of a diffusive field, the contour integral that turns it into
exponential decay in time, and the Gaussian integral that turns it into the
heat kernel in space (\S\ref{sec:hk-fromprop}--\S\ref{sec:hk-kernel}). Then the
two complications a real model can add --- a drift (advection) term
(\S\ref{sec:hk-drift}) and a periodic box (\S\ref{sec:hk-images}). Then the
code that reads the mass, diffusion, and drift symbolically off the inverse
propagator matrix using SageMath (\S\ref{sec:hk-extract}), and the code that
assembles and reconstructs the runtime propagator
(\S\ref{sec:hk-build}). With the propagator in hand we build the tree-level
correlator --- noise driving two retarded legs, collapsed to a single
error-function integral (\S\ref{sec:hk-tree})--- and read the noise strength
out of the action (\S\ref{sec:hk-noise}). Finally we cover what happens at
loop order: the batched analytic transform in the production integrator and the
gate that selects it (\S\ref{sec:hk-loops}), the numerical-transform
cross-check kept alive behind environment variables
(\S\ref{sec:hk-numerical}), the Symanzik-polynomial oracle modules that pin the
production numbers (\S\ref{sec:hk-symanzik}), and the guard rails that make the
whole thing refuse to lie to you (\S\ref{sec:hk-guards}).

\begin{note}
This chapter documents the \emph{real-space propagator layer} for spatial
\msrjd{} theories. The four files at its centre are
\file{msrjd/integration/spatial/heat\_kernel.py} (the closed-form kernels and
the symbolic extraction), \file{spatial\_correlator.py} (the tree-level
correlator and the output transforms), and two modules ---
\file{temporal\_integrate.py} and \file{loop\_parametric.py} --- that carry a
banner reading ``\code{ORACLE-ONLY --- not on the production path}.'' Those two
are independent numerical cross-checks used in tests to pin the production
loop integrator; they are how loops were \emph{validated}, not how loops are
\emph{computed today}. The production loop integrator is
\file{full\_integrator.py}, whose batched transform
\code{\_heat\_kernel\_x\_general} we cover in \S\ref{sec:hk-loops}. We are
careful throughout to keep that distinction visible.
\end{note}

\section{From the inverse propagator to real space}\label{sec:hk-fromprop}

\subsection{The \msrjd{} inverse propagator of a diffusive field}

In the Martin--Siggia--Rose--Janssen--De~Dominicis (\msrjd) formalism a
stochastic PDE is rewritten as a field theory in two fields: the \emph{physical}
field $\phi$ and a \emph{response} (or ``tilde'') field $\tilde\phi$. Take the
prototypical reaction--diffusion (Allen--Cahn-like) equation
\begin{equation}
  \Dt\phi \;=\; B\,\Lap\phi \;-\; A\,\phi \;-\;(\text{nonlinear terms})\;+\;\xi,
  \qquad
  \avg{\xi(x,t)\,\xi(x',t')} \;=\; 2\kappa\,\delta(x-x')\,\delta(t-t').
  \label{eq:hk-langevin}
\end{equation}
The quadratic (``Gaussian'', ``free'') part of the corresponding action gives an
\term{inverse propagator}: the matrix whose inverse is the two-point function.
After a Fourier transform in time and the replacement of the Laplacian operator
by its momentum symbol $\Lap \to -k^2$, the single-field diagonal entry reads
\begin{equation}
  K(\omega,k) \;=\; A \;+\; B\,k^{2} \;+\; \ii\,\omega.
  \label{eq:hk-Kft}
\end{equation}

Three coefficients, three meanings, and the rest of this chapter is in essence
the story of how each one is extracted and used:

\begin{description}
  \item[$A$ --- the \term{mass}.] The $k^0,\omega^0$ term: the relaxation rate
        of the slowest, longest-wavelength mode. For the free theory $A=\mu$.
        At a $\phi^4$ saddle the curvature of the potential shifts it to
        $A=\mu+3\lambda\,\phi_*^{2}$. The code tolerates a \emph{complex} $A$
        (a complex saddle, a Keldysh-rotated combination).
  \item[$B$ --- the \term{diffusion}.] The coefficient of $k^{2}$, i.e.\ the
        $D$ in $D\,\Lap$. Physically it must be positive: a negative $B$ is
        anti-diffusion, which amplifies short wavelengths without bound and is
        ill-posed. The code enforces this (\S\ref{sec:hk-guards}).
  \item[$\ii\,\omega$ --- the time derivative.] The standard \msrjd{}
        first-order-in-time term, and it must be \emph{unit-normalized}: the
        coefficient of $\omega$ has to be exactly $\ii$. The code checks this
        and rejects anything else as ``not a standard \msrjd{} kernel.''
\end{description}

\begin{note}
The textbook \msrjd{} convention sends $\Dt \to -\ii\omega$; this code uses the
$+\ii\omega$ convention, so the time-derivative term in \eqref{eq:hk-Kft}
appears as $+\ii\omega$ rather than $-\ii\omega$. Nothing physical depends on
the sign choice --- it only fixes which half-plane you close the contour in
below --- but it is why the normalization check looks for $\ii$ and not $-\ii$.
\end{note}

The propagator itself is the reciprocal of \eqref{eq:hk-Kft},
\begin{equation}
  G(\omega,k) \;=\; \frac{1}{A + B\,k^{2} + \ii\,\omega}.
  \label{eq:hk-Gft}
\end{equation}

\subsection{Closing the $\omega$-contour: exponential decay in time}

As a function of complex $\omega$, the propagator \eqref{eq:hk-Gft} has a single
simple pole, at $\omega = \ii\,(A + Bk^{2})$. The inverse transform in time,
$G_{\mathrm R}(k,t) = \int \frac{\dd\omega}{2\pi}\,\ee^{-\ii\omega t}\,G(\omega,k)$,
is done by closing the integration contour in the half-plane where the
exponential decays and picking up that one residue. The result is a pure
exponential gated by causality:
\begin{equation}
  \Gret(k,t) \;=\; \theta(t)\,\ee^{-(A + B\,k^{2})\,t}.
  \label{eq:hk-GRkt}
\end{equation}
Here $\theta(t)$ is the Heaviside step --- the retarded (causal) propagator is
identically zero for $t<0$. It is convenient to name the $k$-dependent decay
rate the \term{$k$-dependent mass},
\[
  m(k) \;\equiv\; A + B\,k^{2} \qquad(\text{also written } \mu + D k^{2}),
\]
so that $\Gret(k,t)=\theta(t)\,\ee^{-m(k)t}$.

\section{The heat kernel}\label{sec:hk-kernel}

\subsection{Inverse-transforming $k\to x$}

The remaining transform is from momentum $k$ back to position $x$. The
$k$-dependence of \eqref{eq:hk-GRkt} is the Gaussian $\ee^{-Bk^{2}t}$, and its
inverse spatial Fourier transform is the classic Gaussian-integral identity
\begin{equation}
  \int \frac{\dd^{d}k}{(2\pi)^{d}}\;\ee^{\ii\,k\cdot x}\,\ee^{-B\,k^{2}t}
  \;=\;
  (4\pi B t)^{-d/2}\,\ee^{-|x|^{2}/(4Bt)}.
  \label{eq:hk-gaussianFT}
\end{equation}
This is the \term{heat kernel} --- the Green's function of the diffusion
equation. Multiplying back the $k$-independent decay $\ee^{-At}$ gives the full
real-space retarded propagator:
\begin{equation}
  \boxed{\;
  \Gret(t,x) \;=\; \theta(t)\,(4\pi B t)^{-d/2}\,
  \exp\!\Big[-\frac{|x|^{2}}{4 B t} - A\,t\Big].
  \;}
  \label{eq:hk-GR}
\end{equation}
The entire dependence on the spatial dimension $d$ lives in the prefactor
$(4\pi B t)^{-d/2}$; everything else is dimension-agnostic. Both $A$ and $B$ may
be complex, in which case \eqref{eq:hk-GR} is read as an analytic continuation
of the real formula. This is the single most important equation in the chapter,
and the function that evaluates it is \code{gaussian\_heat\_kernel}.

\begin{defn}
The \term{heat kernel} $(4\pi B t)^{-d/2}\,\ee^{-|x|^{2}/4Bt}$ is the
probability density, at time $t$, of a diffusion that started as a point source
at the origin at $t=0$ with diffusivity $B$. It is the real-space image of a
momentum-space Gaussian. In this subsystem ``analytic IFT'' (inverse Fourier
transform) means exactly this: doing the $k\to x$ transform in closed form via
\eqref{eq:hk-gaussianFT}, rather than sampling the correlator on a grid of
momenta and transforming numerically.
\end{defn}

\subsection{The code: \code{gaussian\_heat\_kernel}}

The closed form \eqref{eq:hk-GR} is implemented essentially verbatim
(\file{heat\_kernel.py:58}). The function takes a scalar time \code{t} (which
may be complex), a scalar separation \code{x} (in $d=1$; for $d\ge 2$ you pass
the radial $|x|$), the mass \code{A} and diffusion \code{B} (both complex-ok),
the dimension \code{spatial\_dim}, and an optional drift \code{V} we return to
in \S\ref{sec:hk-drift}. It returns a Python \code{complex}.

\begin{lstlisting}
def gaussian_heat_kernel(t, x, A, B, spatial_dim: int = 1, V=0.0):
    t = float(t.real) if hasattr(t, 'real') and not isinstance(t, complex) else t
    if (t.real if isinstance(t, complex) else t) <= 0.0:
        return 0j
    A = complex(A); B = complex(B); V = complex(V)
    x_signed = float(x); xx = x_signed ** 2
    pref = (4.0 * math.pi * B * t) ** (-0.5 * spatial_dim)
    drift_exp = 0.0
    if V != 0:
        if spatial_dim != 1:
            raise SpatialPropagatorError(
                'drift heat kernel (V!=0) is implemented for d=1 only ...')
        drift_exp = -1j * V * x_signed / (2.0 * B) + V * V * t / (4.0 * B)
    return complex(pref * cmath.exp(-xx / (4.0 * B * t) - A * t + drift_exp))
\end{lstlisting}

Read it line by line. The first line coerces \code{t} to a real float when it
is a real-like object that is not literally a Python \code{complex} (Sage and
mpmath both produce ``real-like'' objects with a \code{.real} attribute). The
second line is the $\theta(t)$ \term{causality guard}: if the real part of
\code{t} is non-positive the kernel is zero and we return early. The prefactor
line is exactly $(4\pi B t)^{-d/2}$ with $d=$ \code{spatial\_dim}. The final
\code{return} assembles $\text{pref}\cdot\exp(-x^{2}/4Bt - A t + \text{drift})$,
using \code{cmath.exp} (the complex exponential, because $A$ and $B$ can be
complex). The library imports are deliberately minimal --- \code{cmath} and
\code{math} from the standard library (\file{heat\_kernel.py:42--43}) ---
because for the scalar kernel there is nothing exotic to do.

\begin{note}
\textbf{What is \code{cmath}?} It is Python's standard-library module for
complex-valued elementary functions: \code{cmath.exp}, \code{cmath.sqrt}, and so
on. The ordinary \code{math.exp} accepts only real arguments and would raise on
a complex one. Because the mass $A$ and diffusion $B$ here can be complex (a
complex saddle, a Keldysh combination), the kernel must use the complex
exponential. The prefactor $(4\pi B t)^{-d/2}$ is likewise computed with
Python's complex power, so $B<0$ or complex $B$ produces the correct branch
rather than a domain error.
\end{note}

\section{Drift and advection}\label{sec:hk-drift}

A model can carry a first-spatial-derivative term $\partial_x\phi$ ---
advection, a velocity. Under the substitution $\partial_x \to \ii k$ it appears
in the inverse propagator as a term \emph{linear} in $k$, so \eqref{eq:hk-Kft}
generalises to
\begin{equation}
  K(\omega,k) \;=\; A \;+\; V\,k \;+\; B\,k^{2} \;+\; \ii\,\omega,
  \label{eq:hk-Kdrift}
\end{equation}
with $V$ the \term{drift} (the $k^{1}$ coefficient). Completing the square in
the exponent $\ee^{-(A + Vk + Bk^{2})t}$ before doing the Gaussian integral
\eqref{eq:hk-gaussianFT} produces an extra factor in the real-space kernel,
\begin{equation}
  \exp\!\Big[-\frac{\ii V x}{2B} + \frac{V^{2} t}{4B}\Big].
  \label{eq:hk-driftfactor}
\end{equation}
For a drift of the form $V=\ii v$ (the physical advection velocity $v$) this
factor \emph{shifts the centre of the Gaussian to $x = v\,t$} --- the packet is
transported at velocity $v$ while it spreads. In \code{gaussian\_heat\_kernel}
this is the \code{drift\_exp} term, added inside the exponential
(\file{heat\_kernel.py:89}).

Two facts about drift matter for the rest of the manual:

\begin{itemize}
  \item \textbf{$V=0$ is bit-identical to the pure heat kernel.} When the drift
        is zero the \code{if V != 0} branch is skipped entirely and the function
        returns exactly \eqref{eq:hk-GR}. This is not an approximation; it is the
        same arithmetic.
  \item \textbf{Drift is $d=1$ only.} A velocity is a \emph{vector} in $d\ge 2$,
        which is outside the v1 scope, so \code{gaussian\_heat\_kernel} raises a
        \code{SpatialPropagatorError} if it is handed $V\ne 0$ with
        \code{spatial\_dim != 1} (\file{heat\_kernel.py:85--88}).
\end{itemize}

\begin{gotcha}
For a \emph{gradient nonlinearity} --- Burgers $\partial_x(\phi^{2})$, KPZ
$(\partial_x\phi)^{2}$ --- you might expect the propagator to carry a drift. It
usually does not, and that is correct. The only $\partial_x$ that reaches the
bilinear (propagator) sector is the saddle cross-term, which is proportional to
the saddle value $\phi_*$. At the homogeneous saddle $\phi_*=0$ that term
vanishes, so $V\to 0$ and the propagator is the \emph{pure} heat kernel even for
KPZ or Burgers. Do not be surprised to see a ``KPZ'' propagator print
$V(\text{drift})=0$; it is physically right at the trivial saddle.
\end{gotcha}

\section{Periodic boundaries: the method of images}\label{sec:hk-images}

On an infinite line the kernel is \eqref{eq:hk-GR}. On a ring --- a periodic box
of period $L$ --- the Green's function must itself be periodic, and the standard
construction is the \term{method of images}: sum the infinite-domain kernel over
all image sources displaced by integer multiples of $L$,
\begin{equation}
  G_{\text{periodic}}(t,x) \;=\; \sum_{n\in\mathbb{Z}} G_{\infty}(t,\,x + nL).
  \label{eq:hk-imagesum}
\end{equation}
Because the Gaussian tail decays \emph{super-exponentially} in $n$ (the exponent
scales like $-(nL)^{2}$), only a handful of terms matter and the sum truncates
almost immediately. This is the function \code{image\_sum}
(\file{heat\_kernel.py:132}):

\begin{lstlisting}
def image_sum(t, x, A, B, L, spatial_dim=1, eps=1e-12, n_max_cap=2000, V=0.0):
    if spatial_dim != 1:
        raise SpatialPropagatorError(
            'image_sum: periodic BC implemented for d=1 only in v1.')
    base = gaussian_heat_kernel(t, x, A, B, spatial_dim=1, V=V)
    total = base
    ref = abs(base) if abs(base) > 0 else 1.0
    n = 1
    while n <= n_max_cap:
        term_p = gaussian_heat_kernel(t, x + n * L, A, B, spatial_dim=1, V=V)
        term_m = gaussian_heat_kernel(t, x - n * L, A, B, spatial_dim=1, V=V)
        total += term_p + term_m
        if abs(term_p) < eps * ref and abs(term_m) < eps * ref and n >= 2:
            break
        n += 1
    return complex(total)
\end{lstlisting}

The $n=0$ term \code{base} sets a reference magnitude \code{ref}; the loop adds
the $\pm n$ images in pairs and stops once both have fallen below
\code{eps}~$\times$~\code{ref} (default $10^{-12}$) \emph{and} at least two
shells have been added (the \code{n >= 2} guard avoids a premature stop at a
node where a single term happens to be tiny). A hard cap \code{n\_max\_cap}
prevents a runaway. Like drift, periodic boundaries are $d=1$ only here:
$d\ge 2$ periodic transforms would need a lattice sum per axis, and the function
raises otherwise (\file{heat\_kernel.py:143--145}). The $d\ge 2$ periodic case
is handled instead by a separate lattice-sum transform, \code{periodic\_inverse\_ft}, in \S\ref{sec:hk-numerical}.

\section{Reading $(A,B,V)$ off the inverse propagator with Sage}\label{sec:hk-extract}

So far we have assumed the three coefficients $(A,B,V)$ are known. In practice
they arrive as \emph{symbolic} expressions buried inside the inverse-propagator
matrix \code{K\_ft} that the shared \code{build\_propagator} produces upstream:
a single diagonal entry might look like
\code{mu + 3*lambda*phistar1\^{}2 - D*Laplacian + I*omega}. Peeling
$(A,B,V)$ out of that is a small computer-algebra task, and the tool for it is
SageMath.

\begin{note}
\textbf{What is SageMath?} Sage is a large open-source mathematics system built
on top of Python; it bundles dozens of mathematical libraries behind one Python
API. The only piece this subsystem uses is Sage's \term{Symbolic Ring}, written
\code{SR} --- a computer-algebra engine for symbolic expressions, in the spirit
of Mathematica but callable from Python. A Sage symbolic expression can hold
unknowns (\code{mu}, \code{D}, \code{phistar1}, \code{k}, \code{omega}), be
expanded and substituted into, and have its polynomial coefficients read off.
The imports are \code{from sage.all import SR, I as SR\_I, matrix}
(\file{heat\_kernel.py:47}): \code{SR} is the ring itself, \code{SR\_I} is the
symbolic imaginary unit $\ii$ (renamed so it does not clash with loop counters
named \code{i} or \code{I}), and \code{matrix} builds a matrix whose entries
live in \code{SR}. The reason this layer stays in Sage rather than switching to
sympy is that \code{build\_propagator} already hands it a Sage matrix, so
staying in Sage avoids a translation round-trip.
\end{note}

\subsection{The single-entry extractor}

The workhorse is \code{extract\_mass\_diffusion} (\file{heat\_kernel.py:161}),
which takes one diagonal entry of \code{K\_ft} together with the relevant
symbols and returns the triple $(A,B,V)$. The substitution and guards are worth
quoting:

\begin{lstlisting}
subs = {lap_sym: -k_var**2}
if grad_sym is not None:
    subs[grad_sym] = SR_I * k_var          # d_x -> i k  (the drift symbol)
e = SR(kft_entry).subs(subs).expand()
if e.has(lap_sym) or (grad_sym is not None and e.has(grad_sym)):
    raise SpatialPropagatorError(...)      # a residual operator symbol survived

omega_coeff = e.coefficient(omega, 1)      # the coefficient of omega^1
if (omega_coeff - SR_I).simplify_full() != 0:
    raise SpatialPropagatorError(...)      # not +i*omega-normalized

e0 = e.coefficient(omega, 0)               # the omega-independent part
deg = e0.degree(k_var)
if deg > 2:
    raise SpatialPropagatorError(...)      # higher-derivative operator
B = e0.coefficient(k_var, 2)               # k^2 coefficient = diffusion
V = e0.coefficient(k_var, 1)               # k^1 coefficient = drift
A = e0.coefficient(k_var, 0)               # k^0 coefficient = mass
return A, B, V
\end{lstlisting}

Each Sage call earns its keep:

\begin{description}
  \item[\code{.subs(dict)}] substitutes one expression for another. Here the
        inert Laplacian symbol becomes $-k^{2}$ and, if a first-derivative symbol
        \code{grad\_sym} is present, it becomes $\ii k$. After the substitution
        the entry is an ordinary polynomial in $k$ and $\omega$.
  \item[\code{.expand()}] multiplies everything out so the coefficients can be
        read individually.
  \item[\code{.has(sym)}] is \code{True} iff the expression still contains
        \code{sym}. If a Laplacian (or gradient) symbol survives the
        substitution, the entry encoded a higher-derivative operator (a
        $(\Lap)^{2}$ would leave a $k^{4}$ and a residual symbol), and the function
        rejects it.
  \item[\code{.coefficient(var, n)}] returns the coefficient of $\code{var}^{n}$.
        This is how $A,B,V$ and the $\ii\omega$ coefficient are picked off.
  \item[\code{.simplify\_full()}] runs Sage's strongest simplifier so that the
        normalization test $(\text{omega\_coeff}-\ii)$ genuinely collapses to
        \code{0} when it should. A naive subtraction might leave an
        unsimplified expression that is nonzero only syntactically.
  \item[\code{.degree(k\_var)}] returns the highest power of $k$. The guard
        \code{deg > 2} is what enforces ``no higher-derivative operators'' (see
        \S\ref{sec:hk-guards}).
\end{description}

The mapping from these three coefficients to physics is exactly the dictionary
of \S\ref{sec:hk-fromprop}: \code{B} is the $k^{2}$ coefficient (diffusion),
\code{V} the $k^{1}$ coefficient (drift), \code{A} the $k^{0}$ coefficient
(mass). For the worked example above, $A=\code{mu + 3*lambda*phistar1\^{}2}$,
$B=\code{D}$, and $V=0$.

\begin{gotcha}
The gradient symbol is fetched as \code{grad\_sym = SR.var('GradX')}
(\file{heat\_kernel.py:360}). Sage \emph{caches symbolic variables by name}: any
two calls to \code{SR.var('GradX')} anywhere in the codebase return the
\emph{same} object. This is load-bearing --- the gradient symbol declared when
the operator IR is lowered and the one declared here must be identical, or the
\code{.subs} would silently fail to match. If you ever rename it in one place
you must rename it everywhere.
\end{gotcha}

\subsection{The coupled-field generalization}

\code{extract\_mass\_diffusion} reads \emph{one} diagonal entry. Its sibling
\code{reaction\_diffusion\_matrices} (\file{heat\_kernel.py:226}) does the same
for the \emph{full} matrix, returning $N\times N$ Sage matrices $(M,\mathcal{D},V)$
of masses, diffusions, and drifts. The extra requirement at the matrix level is
that the time derivative be \emph{diagonal and unit-normalized}: the $\ii\omega$
coefficient must be $\ii$ on the diagonal and exactly $0$ off it
(\file{heat\_kernel.py:262--268}). For a diagonal \code{K\_ft} this reduces, per
entry, to \code{extract\_mass\_diffusion} bit-identically. These matrices feed
the coupled-field spectral propagator (the subject of
Chapter~\ref{ch:spatial-coupled}); the diagonal Tier-1 path we are documenting
here does \emph{not} consume them.

\section{From propagator to runtime: build versus reconstruct}\label{sec:hk-build}

\subsection{Building the picklable spatial block}

\code{build\_spatial\_propagator} (\file{heat\_kernel.py:310}) is the assembler.
It is handed the inverse-propagator matrix \code{K\_ft}, the $\omega$ symbol, the
namespace \code{ns} that carries \code{ns.Laplacian}, the model dict, and the
field-name lists. It reads the spatial dimension, declares the $k$ and
\code{GradX} symbols, runs the diagonality and boundary guards, calls
\code{extract\_mass\_diffusion} on each diagonal entry, and returns a plain
dictionary --- the \emph{spatial block} of the propagator dict. The single most
important property of that dictionary is that it is \textbf{picklable}: it holds
only Sage expressions and plain data, no closures, so it caches cleanly through
the pipeline's \code{PipelineCache}.

The block it returns (\file{heat\_kernel.py:421--431}) has these keys:

\begin{center}
\begin{longtable}{@{}p{0.27\textwidth} p{0.66\textwidth}@{}}
\toprule
\textbf{key} & \textbf{meaning} \\
\midrule
\code{G\_tx\_sym} & \code{dict[(i,j)] -> (A\_expr, B\_expr)} or \code{None};
  per-entry symbolic mass/diffusion; \code{None} off-diagonal. \\
\code{k\_var} & the Sage momentum symbol $k$. \\
\code{spatial\_dim} & the dimension $d\in\{1,2,3\}$. \\
\code{bc\_mode} & \code{'infinite'} or \code{'periodic'}. \\
\code{bc\_params} & boundary parameters (e.g.\ \code{\{'length': 'L'\}}),
  with \code{'mode'} stripped out. \\
\code{initial\_mode} & \code{'stationary'} (default), etc. \\
\code{ac\_mass} & \code{dict[i] -> A\_expr}; per-field mass. \\
\code{ac\_diffusion} & \code{dict[i] -> B\_expr}; per-field diffusion. \\
\code{ac\_drift} & \code{dict[i] -> V\_expr}; per-field drift ($0$ means a pure
  heat kernel). \\
\bottomrule
\end{longtable}
\end{center}

Note that \code{G\_tx\_sym} keeps each diagonal entry as a 2-tuple $(A,B)$ for
backward compatibility; the drift $V$ travels separately in \code{ac\_drift}.
The off-diagonal entries of \code{G\_tx\_sym} are set to \code{None}, and the
diagonality guard (\file{heat\_kernel.py:369--377}) guarantees there is nothing
there to lose --- any nonzero off-diagonal \code{K\_ft[i,j]} raises a
\code{SpatialPropagatorError} before we ever get here, because the closed-form
Tier-1 kernel requires each field to be its own independent Allen--Cahn mode.

\subsection{Why the callables are not stored}

The one thing the block does \emph{not} contain is the runtime propagator
function itself --- the thing you actually call as $G(t,x)$. The reason is
purely mechanical: \textbf{closures do not pickle}. A function that has captured
\code{A\_num}, \code{B\_num}, and the boundary mode in its enclosing scope cannot
be serialized, so it cannot live in a cacheable dict. The fix is to store only
the symbolic data (which \emph{does} pickle) and rebuild the callables on demand.

That rebuild is \code{make\_g\_tx\_callables} (\file{heat\_kernel.py:434}). It is
idempotent and cheap, and it must be called after \code{build\_propagator}
returns --- whether that was a fresh build or a cache load. It walks
\code{G\_tx\_sym}, turns each symbolic $(A,B)$ (and the matching drift $V$) into
numeric callables via the helper \code{\_make\_numeric}, and wraps them in a
per-entry closure that dispatches on the boundary mode:

\begin{lstlisting}
def _g(t, x, **num_params):
    A = A_num(**num_params)
    B = B_num(**num_params)
    V = V_num(**num_params) if V_num is not None else 0.0
    if bc_mode == 'periodic':
        ...
        L = (num_params[L_name] if isinstance(L_name, str) else float(L_name))
        return image_sum(t, x, A, B, float(L), spatial_dim=d, V=V)
    return gaussian_heat_kernel(t, x, A, B, spatial_dim=d, V=V)
\end{lstlisting}

The closure evaluates the symbolic coefficients at the supplied numeric
parameters and then calls \code{image\_sum} (periodic) or
\code{gaussian\_heat\_kernel} (infinite). The helper \code{\_make\_numeric}
(\file{heat\_kernel.py:287}) is the bridge from symbol to number: it reads an
expression's free symbols with \code{expr.variables()}, sorts them by name, and
returns a closure that substitutes each by name from a \code{num\_params} dict
and coerces the result to a Python \code{complex} --- raising a \code{KeyError}
with a helpful message if a parameter or saddle value is missing. Entries whose
symbol is \code{None} (the off-diagonal ones) get a closure that simply returns
\code{0j}.

\begin{example}\label{ex:hk-gtx}
A concrete round trip. After building a $\phi^4$ propagator at the trivial
saddle and reconstructing the callables, calling the diagonal entry as
\begin{center}
\code{G\_tx[(0,0)](t=0.5, x=1.0, mu=1.0, lambda=0.0, phistar1=0.0, D=0.1)}
\end{center}
substitutes $A=\mu=1$, $B=D=0.1$, $V=0$ into \eqref{eq:hk-GR} and returns
$(4\pi\cdot 0.1\cdot 0.5)^{-1/2}\,\exp[-1/(4\cdot 0.1\cdot 0.5) - 0.5]$.
\end{example}

\section{The tree-level correlator}\label{sec:hk-tree}

\subsection{Noise driving two retarded legs}

The retarded propagator \eqref{eq:hk-GR} is the response to a kick. The
\emph{correlation} function $C(x,\tau)=\avg{\phi(x_1,t_1)\,\phi(x_2,t_2)}$ (with
$x=x_1-x_2$, $\tau=t_1-t_2$) is, at tree level, the \emph{noise sourcing two
retarded propagators}:
\begin{equation}
  C_{ij}(x,\tau) \;=\; 2\,D^{\text{noise}}_{ij}
  \int_{-\infty}^{\min(t_1,t_2)}\!\dd t_v \int \dd x_v\;
  {\Gret}_i(t_1-t_v,\,x-x_v)\,{\Gret}_j(t_2-t_v,\,-x_v).
  \label{eq:hk-Ctwoleg}
\end{equation}
The noise, injected at the vertex $(x_v,t_v)$ and white in space and time,
propagates causally to both external points. For a diagonal heat-kernel
propagator the spatial integral over $x_v$ collapses by the heat-kernel
\term{semigroup} property --- the convolution of two Gaussians is one Gaussian
--- and the remaining $t_v$ integral reduces to a single one-dimensional
integral of error-function type,
\begin{equation}
  C_{ii}(x,\tau) \;=\; \frac{D^{\text{noise}}_i}{\sqrt{4\pi B_i}}
  \int_{|\tau|}^{\infty} s^{-1/2}\,
  \exp\!\Big[-\frac{x^{2}}{4 B_i s} - A_i s\Big]\,\dd s.
  \label{eq:hk-Cerf}
\end{equation}

\subsection{The error-function closed form}

The integral $\int s^{-1/2}\,\ee^{-\beta/s - \alpha s}\,\dd s$ in
\eqref{eq:hk-Cerf} has a genuine closed form. Substituting $s=w^{2}$, the
antiderivative of $\ee^{-\alpha w^{2} - \beta/w^{2}}$ is the ``erf-split'' form
\begin{equation}
  F(w) \;=\; \frac{\sqrt\pi}{4\sqrt\alpha}\Big[
    \ee^{+2\sqrt{\alpha\beta}}\,\mathrm{erf}\!\big(\sqrt\alpha\,w + \tfrac{\sqrt\beta}{w}\big)
    + \ee^{-2\sqrt{\alpha\beta}}\,\mathrm{erf}\!\big(\sqrt\alpha\,w - \tfrac{\sqrt\beta}{w}\big)
  \Big],
  \label{eq:hk-Fw}
\end{equation}
so the definite integral equals $2\,[F(\sqrt{U})-F(\sqrt{L})]$, with
$\alpha=$ mass, $\beta=x^{2}/4B\ge 0$. The semi-infinite case $U\to\infty$
(valid when $\operatorname{Re}\sqrt\alpha>0$) has the simple limit
$F(\infty)=\tfrac{\sqrt\pi}{4\sqrt\alpha}\big(\ee^{2\sqrt{\alpha\beta}}+\ee^{-2\sqrt{\alpha\beta}}\big)$.
This is the function \code{erf\_time\_integral} (\file{heat\_kernel.py:93}).

\begin{gotcha}
The form \eqref{eq:hk-Fw} contains a \emph{near-cancellation of two large
exponentials}: when $\sqrt{\alpha\beta}$ is large, $\ee^{+2\sqrt{\alpha\beta}}$
and $\ee^{-2\sqrt{\alpha\beta}}$ differ by many orders of magnitude, and double
precision loses most of its significant digits in the subtraction. For that
reason \code{erf\_time\_integral} is evaluated at \textbf{30 decimal digits}
using \code{mpmath} and only then collapsed back to a \code{complex}. If you
ever port this to plain floats, expect it to silently lose precision.
\end{gotcha}

\begin{note}
\textbf{What is \code{mpmath}?} It is a pure-Python \emph{arbitrary-precision}
floating-point library: where ordinary floats give about 16 significant digits,
mpmath gives 30, 50, or 1000 on demand, and its special functions (\code{erf},
\code{sqrt}, \code{exp}) work on \emph{complex} arguments. Sage ships it. The
import is \code{import mpmath as mp} (\file{heat\_kernel.py:46}), and the whole
high-precision computation runs inside a \code{with mp.workdps(dps):} block ---
\code{workdps} is a context manager that temporarily sets the working precision
to \code{dps} decimal digits. The complex numbers are \code{mp.mpc(...)}, reals
are \code{mp.mpf(...)}.
\end{note}

The implementation also handles the $w\to 0^{+}$ endpoint carefully. The helper
\code{\_F(w)} (\file{heat\_kernel.py:115--122}) special-cases \code{w == 0}:
$\mathrm{erf}(\pm\infty)=\pm 1$ when $\beta>0$, and $\mathrm{erf}(0)=0$ when
$\beta=0$. This is the $s\to 0$ endpoint that appears whenever the lower limit
$L$ is zero (the equal-time, $x=0$ corner), and getting it wrong would produce a
spurious $0/0$ right at the most physically interesting point of the correlator.

\subsection{\code{free\_two\_point}: one field, both boundary conditions}

\code{free\_two\_point} (\file{spatial\_correlator.py:138}) wraps the erf
integral into a per-field tree correlator. Its signature is
\code{free\_two\_point(mu, D, kap, x, tau, bc\_mode='infinite', L=None, ...)},
where \code{mu} is the mass, \code{D} the diffusion, and \code{kap} the noise
spectral coefficient:

\begin{lstlisting}
def free_two_point(mu, D, kap, x, tau, bc_mode='infinite', L=None, n_images_cap=200):
    D = float(D); mu = complex(mu); kapf = float(kap)
    pref = kapf / math.sqrt(4.0 * math.pi * D)
    def _C_inf(xx):
        beta = (xx * xx) / (4.0 * D)
        integ = erf_time_integral(mu, beta, abs(tau), U_hi=None)
        return complex(pref * integ)
    if bc_mode == 'infinite':
        return _C_inf(x)
    if bc_mode == 'periodic':
        ...                                # sum images _C_inf(x +/- nL)
\end{lstlisting}

For the infinite domain it returns the single erf integral with
$\beta=x^{2}/4D$ and lower limit $|\tau|$ (the $U\to\infty$ semi-infinite form,
\code{U\_hi=None}). For a periodic box it sums the erf correlator over images
$x\to x\pm nL$ with the same super-exponential truncation as \code{image\_sum}.

\begin{gotcha}
The first two arguments of \code{free\_two\_point} are named \code{(mu, D)}, but
its callers pass them the mass and diffusion \emph{they} call $(A,B)$. In
\code{compute\_spatial\_correlator\_tree} the call is
\code{free\_two\_point(A, B, Dn, ...)} (\file{spatial\_correlator.py:291}). This
is semantically correct --- mass maps to \code{mu}, diffusion maps to \code{D}
--- but the name change can read as a bug at a glance. The rule to remember:
\emph{the first argument is always the mass, the second is always the
diffusion}, whatever the local letter.
\end{gotcha}

\subsection{Assembling the $(x,\tau)$ grid}

\code{compute\_spatial\_correlator\_tree} (\file{spatial\_correlator.py:178}) is
the driver that fills the output array. It takes the expanded FieldTheory, the
model, the propagator dict, the numeric parameters, the two external legs, and
the $(\tau,x)$ grids, and returns a complex array of shape
\code{(len(tau\_grid), len(spatial\_grid))} together with an \code{info} dict.
Its logic, in order:

\begin{enumerate}
  \item Require a spatial block (\code{prop['spatial\_dim']}); normalize the
        numeric parameters to string keys; resolve the periodic length $L$ by
        name or inline value.
  \item \textbf{Resolve the field index} of the external legs by matching the
        leg base name to the physical column names (stripping trailing
        population digits).
  \item \textbf{Tier-1 guard:} if \code{prop['G\_tx\_sym'] is None} the theory
        is not the diagonal heat-kernel case, and the function raises a clear
        \code{NotImplementedError} explaining that off-diagonal coupling and
        higher-derivative operators are v2 features
        (\file{spatial\_correlator.py:239--248}).
  \item Read $(A,B)$ for the resolved diagonal entry, substitute parameters, and
        take $A$ complex, $B$ real.
  \item \textbf{Diffusion-sign guard:} raise on $B<0$ (anti-diffusive, the heat
        kernel diverges) and on $B=0$ (no Laplacian: a time-only field, for
        which a spatial correlator is undefined). See \S\ref{sec:hk-guards}.
  \item Read the noise coefficient $D_{\text{noise}}$ for this field
        (\S\ref{sec:hk-noise}); raise if absent.
  \item Fill \code{C[it, ix] = free\_two\_point(A, B, Dn, x, tau, ...)} over the
        grid.
\end{enumerate}

The static ($\tau=0$) limit of \eqref{eq:hk-Cerf} has its own closed form that
the erf result must reproduce --- $C(x,0)=\frac{D^{\text{noise}}}{2\sqrt{AB}}\,
\ee^{-|x|\sqrt{A/B}}$ in $d=1$ --- and the function
\code{free\_correlator\_static\_closed\_form} (\file{spatial\_correlator.py:398})
supplies it in all three dimensions: the $d=1$ exponential, the $d=2$ modified
Bessel $K_0$, and the $d=3$ Yukawa $\ee^{-\kappa r}/r$. It serves as an
independent oracle for the transforms of \S\ref{sec:hk-numerical}.

\section{Reading the noise strength out of the action}\label{sec:hk-noise}

The correlator \eqref{eq:hk-Cerf} needs the noise strength $D^{\text{noise}}$,
and that lives in the action, not in the propagator. In the \msrjd{} action a
white-noise covariance $\avg{\xi\xi}=2\kappa\,\delta$ appears as a term
$-D^{\text{noise}}\,\tilde\phi^{2}$ --- a term with \emph{two response fields and
zero physical fields}. The framework classifies every action term by its
``bigrade'' $(\text{number of response factors},\,\text{number of physical
factors})$, so the noise sector is precisely the $(2,0)$ bigrade.

\code{extract\_noise\_coefficients} (\file{spatial\_correlator.py:46}) reads it.
It fetches the $(2,0)$ sector from the expanded FieldTheory
(\code{ft.\_by\_tp[(2,0)]}), iterates the monomials of that polynomial, and
identifies a \emph{pure} $\tilde\phi_i^{2}$ monomial as one with exponent $2$ on
ring position $i$ (response fields are the first \code{n\_tilde} generators of
the ring) and $0$ everywhere else. For such a monomial the white-noise
coefficient is $D^{\text{noise}}_i = -\operatorname{Re}(\text{coeff})$ --- the
minus sign because the action term is $-D^{\text{noise}}\tilde\phi^{2}$. The
return value is a dict \code{\{field\_index: D\_noise\}}. This convention is
exactly the one that reproduces the scalar OU result $\avg{x^{2}}=D/\mu$ from an
action term $-D\,\tilde x^{2}$.

Its coupled-field sibling \code{extract\_noise\_matrix}
(\file{spatial\_correlator.py:93}) returns the full $n_{\tilde\phi}\times
n_{\tilde\phi}$ noise-covariance matrix $N$ with the convention
$N_{ii}=-2\,\mathrm{coeff}(\tilde\phi_i^{2})$ and
$N_{ij}=-\mathrm{coeff}(\tilde\phi_i\tilde\phi_j)$ for $i\ne j$, so that
$N_{ii}=2\kappa_i$ with $\kappa_i$ the value the diagonal extractor returns.
Both extractors read only the $q^{0}$ (white-noise) part; $q^{2}$-dependent
(conserved, Model-B) noise is skipped and left at zero, matching the diagonal
heat-kernel scope.

\section{The analytic IFT at loop order}\label{sec:hk-loops}

Everything to here has been tree level. At loop order the correlator picks up a
self-energy correction $\delta C$, and the central question becomes how to do the
spatial transform of a \emph{loop} diagram. The answer is the same trick at a
larger scale: do the loop momentum integral analytically (it is Gaussian), do
the external $q\to x$ transform analytically (it is again Gaussian), and never
build a $q$-grid.

\subsection{The per-chamber Gaussian}

A loop diagram, after its loop momentum integral is done in Schwinger
parameters (the machinery of \S\ref{sec:hk-symanzik} and the spatial-core
chapter), leaves a residual Gaussian in the \emph{external} momentum $q$ of the
form $\ee^{-q^{\mathsf T}\mathcal{B}\,q}$, where for each Schwinger sample
$\mathcal{B}=D\cdot Q_{\text{eff}}$ (a scalar for a two-point function, an
$n\times n$ matrix for $k\ge 3$ external legs). Inverse-transforming that
Gaussian to real space is once more the heat-kernel identity
\eqref{eq:hk-gaussianFT}. The production code does this for a whole \emph{batch}
of Schwinger samples at once, vectorized, in \code{\_heat\_kernel\_x\_general}
(\file{full\_integrator.py:143}):

\begin{lstlisting}
def _heat_kernel_x_general(Bcal_g, xs_arr, spatial_dim):
    d = float(spatial_dim)
    xs_arr = np.asarray(xs_arr, dtype=float)
    if Bcal_g.ndim == 1:                                     # n_ext = 1 (k=2)
        if xs_arr.ndim == 1:
            x2 = xs_arr ** 2
        else:
            x2 = np.sum(xs_arr ** 2, axis=-1)
        return ((4.0 * math.pi * Bcal_g)[:, None] ** (-0.5 * d)
                * np.exp(-x2[None, :] / (4.0 * Bcal_g[:, None])))
    n = Bcal_g.shape[1]
    ...
    detB = np.linalg.det(Bcal_g)                             # (P',)
    Binv = np.linalg.inv(Bcal_g)                             # (P', n, n)
    if xs_arr.ndim == 2:                                     # d=1: (n_x, n)
        quad = np.einsum('xj,pjk,xk->px', xs_arr, Binv, xs_arr)
    else:                                                    # d>=2: (n_x, n, d)
        quad = np.einsum('xjc,pjk,xkc->px', xs_arr, Binv, xs_arr)
    return ((4.0 * math.pi) ** (-0.5 * n * d)
            * detB[:, None] ** (-0.5 * d)
            * np.exp(-0.25 * quad))
\end{lstlisting}

For a two-point function (\code{n\_ext = 1}) this is exactly the body of
\code{gaussian\_heat\_kernel} --- $(4\pi B)^{-d/2}\,\ee^{-|x|^{2}/4B}$ --- with
the $\theta(t)$, drift, and time pieces already folded into $\mathcal{B}$
upstream, and with the batch axis $P'$ over Schwinger samples vectorized. For
$k\ge 3$ external legs (\code{n\_ext >= 2}) it is the \emph{multivariate}
Gaussian inverse transform
\[
  \text{hk} \;=\; (4\pi)^{-nd/2}\,\det(\mathcal{B})^{-d/2}\,
  \exp\!\Big[-\tfrac14\sum_c \vec X_c^{\mathsf T}\,\mathcal{B}^{-1}\,\vec X_c\Big],
\]
one factor per spatial component $c$ (the components share the same
$\mathcal{B}$ by isotropy). The determinant and inverse are batched with
\code{np.linalg.det} / \code{np.linalg.inv}, and the quadratic form is contracted
with \code{np.einsum}.

\begin{note}
\textbf{What is \code{np.einsum}?} It is NumPy's Einstein-summation routine:
you write the index pattern of a tensor contraction as a string and it performs
exactly that sum. Here \code{'xj,pjk,xk->px'} reads ``contract the evaluation
point $x$'s component $j$ with $\mathcal{B}^{-1}_{jk}$ and the component $k$,
summing over $j$ and $k$, for every sample $p$,'' producing the array of
quadratic forms $\vec X^{\mathsf T}\mathcal{B}^{-1}\vec X$ indexed by sample $p$
and point $x$. The \code{[:, None]} index inserts a length-one axis so the
per-sample prefactor broadcasts against the per-point exponential.
\end{note}

\subsection{The gate: \code{\_use\_analytic}}

Whether the analytic transform is used at all is decided by one line in the
orchestrator \file{pipeline\_bridge.py}. The relevant flags
(\file{pipeline\_bridge.py:1667,1677}) are:

\begin{lstlisting}[style=console]
_all_plain    = all(rec[2] is None for rec in live_g)   # no derivative vertices
_force_num    = _os.environ.get('SPATIAL_FORCE_NUMERICAL_FT', '') == '1'
_use_analytic = (_all_plain or d == 1) and not _force_num
\end{lstlisting}

So the analytic heat-kernel transform covers two regimes: \textbf{(a)} plain
(derivative-free) vertices at any dimension $d$, and \textbf{(b)} any vertices at
$d=1$, where the derivative-vertex form factors are handled by a polynomial fit
plus closed-form heat-kernel $q$-moments. The one remaining gap is $d\ge 2$
\emph{with} derivative vertices, which still falls back to the numerical
transform of \S\ref{sec:hk-numerical}. When \code{\_use\_analytic} is true, each
diagram's per-chamber Schwinger samples produce $\mathcal{B}=D\cdot Q_{\text{eff}}$,
which is fed straight to \code{\_heat\_kernel\_x\_general} to get $\delta C(x,\tau)$
directly --- with no $q$-grid anywhere.

\section{The numerical-transform cross-check and the env knobs}\label{sec:hk-numerical}

The analytic transform is exact, but ``exact'' is a claim that has to be
\emph{checked}, and the check is a second, independent route: do the $q\to x$
transform numerically on a discretized momentum grid. That route lives in
\code{radial\_inverse\_ft} and \code{periodic\_inverse\_ft}
(\file{spatial\_correlator.py:299,341}).

\code{radial\_inverse\_ft} takes an isotropic momentum-space correlator
$C(|q|)$ sampled on a $q$-grid and transforms it to $C(|x|)$ via the radial
inverse transform appropriate to the dimension:
\begin{align*}
  d=1:&\quad C(x) = \tfrac{1}{\pi}\int_0^{\infty}\cos(qx)\,C(q)\,\dd q,\\
  d=2:&\quad C(x) = \tfrac{1}{2\pi}\int_0^{\infty} q\,J_0(qx)\,C(q)\,\dd q
       \quad(\text{Hankel order }0),\\
  d=3:&\quad C(x) = \tfrac{1}{2\pi^{2}x}\int_0^{\infty} q\,\sin(qx)\,C(q)\,\dd q,
\end{align*}
each evaluated as a \code{np.trapz} (trapezoidal numerical integration) over the
grid, with the $d=3$ $x\to 0$ limit ($\sin(qx)/x\to q$) handled explicitly. The
$d=2$ branch uses the Bessel function $J_0$ from \code{scipy.special.j0}.
\code{periodic\_inverse\_ft} is the discrete-momentum analogue for a periodic
cubic box: it builds the integer-mode lattice with \code{np.meshgrid},
interpolates the radial $C(q)$ onto the lattice magnitudes with \code{np.interp},
zeros out modes past the cutoff with \code{np.where}, and sums
$\cos(k_x x)\,C(|k_n|)$ with a $1/L^{d}$ prefactor. The two boundary conditions
share one momentum-space correlator and differ only in the transform.

\begin{gotcha}
Both numerical transforms \textbf{truncate at $k_{\max}=q\_grid[-1]$} --- they
evaluate the correlator \emph{at} a finite momentum cutoff, not in the
continuum. The continuum value is recovered only in the $k_{\max}\to\infty$,
fine-grid limit. This is the ``Regime 1'' (physical-cutoff) interpretation. The
\emph{analytic} heat-kernel transform has no such truncation --- that is the
whole point of ``retiring the $q$-grid'' --- so the two routes agree only as the
numerical grid is refined.
\end{gotcha}

Three environment variables, all read in \file{pipeline\_bridge.py}, control the
cross-check:

\begin{description}
  \item[\code{SPATIAL\_FORCE\_NUMERICAL\_FT}] (\file{pipeline\_bridge.py:610,1676})
        forces the numerical-transform path even where the analytic transform
        would apply (it is the \code{not \_force\_num} term in the gate above).
        This keeps the cosine/radial cross-check reachable. It is exact only in
        the $q\_cut\to\infty$, $n\_q\to\infty$ limit; use it for validation, not
        production.
  \item[\code{SPATIAL\_Q\_CUT}] (\file{pipeline\_bridge.py:1604}) overrides the
        $q$-grid upper cutoff for the numerical cross-check. The in-code comment
        warns that derivative-vertex form factors with a $q^{4}$ tail need a
        large cutoff for the truncated trapezoid to converge --- a limit the
        analytic path does not have.
  \item[\code{SPATIAL\_N\_Q}] (\file{pipeline\_bridge.py:1605}) overrides the
        number of $q$-grid points for the cross-check.
\end{description}

\begin{gotcha}
A separate guard on the same code path can refuse to run at all. Before
allocating a chamber's $n_t^{\,n_V}\cdot n_s^{\,n_C}\cdot n_x$ complex array,
\file{pipeline\_bridge.py:1641--1665} checks it against
\code{SPATIAL\_MEM\_BUDGET\_GB} (default 6\,GB) and aborts up front if it would
exceed the budget. This is a deliberate guard against an out-of-memory event
that can crash not just the kernel but the OS --- the same macOS hazard that
forced the spatial path to be thread-based rather than fork-based. It is not in
the four primary files of this chapter, but it governs whether the analytic
transform even gets to run.
\end{gotcha}

\section{The Symanzik momentum core (oracle path)}\label{sec:hk-symanzik}

We now reach the two modules that carry the \code{ORACLE-ONLY} banner. They
implement the loop \emph{momentum} integral in Schwinger parameters as an
independent numerical reference, and they exist to pin the production
\code{full\_integrator} numbers in tests. \emph{They are not on the production
path}: \code{compute\_cumulants} does not call them.

\subsection{Why momentum-first}

The motivation (stated in the module docstring,
\file{loop\_parametric.py:8--31}) is precision. Doing the loop integral
time-first reuses the temporal Phase~J machinery plus a numerical $\int\dd\ell$,
but spatial loops \emph{generically} trip the $m\ge 3$ close-pair precision slow
path, because the loop momentum sweeps the edge masses past one another. The cure
is to do the momentum integral \emph{analytically first}. With a Schwinger
parameter $w_e$ on each edge, every edge is $\ee^{-(\mu + D k_e^{2})w_e}$, and if
each edge momentum is routed as $k_e = a_e\,\ell + b_e\,q$ (one loop momentum
$\ell$, external $q$), the loop integral is a pure Gaussian:
\begin{equation}
  \int \frac{\dd^{d}\ell}{(2\pi)^{d}}\;\ee^{-D\sum_e w_e k_e^{2}}
  \;=\;
  (4\pi D U)^{-d/2}\,\ee^{-D q^{2}(W - V^{2}/U)},
  \label{eq:hk-symanzik1}
\end{equation}
with the \term{Symanzik forms}
\[
  U = \sum_e a_e^{2} w_e,\qquad
  V = \sum_e a_e b_e w_e,\qquad
  W = \sum_e b_e^{2} w_e,\qquad
  F_{\text{reduced}} = W - V^{2}/U.
\]
There are \emph{no momentum-dependent poles} here --- only a smooth
Schwinger/time integral remains, so close-pair singularities can never arise.
This is the single conceptual reason the spatial loop integrator works at all.

\subsection{\code{loop\_parametric.py} and \code{spatial\_reduce.py}}

\code{loop\_parametric.symanzik\_UF} (\file{loop\_parametric.py:41}) is the
$L=1$ core: it computes $U,V,W,F_{\text{reduced}}$ and the prefactor
$(4\pi D U)^{-d/2}$, raising on $U\le 0$ (which would mean no loop-momentum
damping). \code{gaussian\_momentum\_integral} (\file{loop\_parametric.py:73})
is the thin wrapper $\text{pref}\cdot\ee^{-Dq^{2}F_{\text{reduced}}}$. The
$L$-loop generalization lives one module over in \code{spatial\_reduce.py},
which replaces the scalars by matrices --- $\Lambda$ $(L\times L)$, $N$
$(L\times n_{\text{ext}})$, $Q$ $(n_{\text{ext}}\times n_{\text{ext}})$ --- with
$U=\det\Lambda$ and the reduced external form
$Q_{\text{eff}}=Q-N^{\mathsf T}\Lambda^{-1}N$. Its \code{momentum\_integral}
(\file{spatial\_reduce.py:101}) evaluates
\[
  \int \prod_i\frac{\dd^{d}\ell_i}{(2\pi)^{d}}\;\ee^{-D\sum_e w_e k_e^{2}}
  = (4\pi D)^{-Ld/2}\,U^{-d/2}\,\ee^{-D q^{\mathsf T} Q_{\text{eff}} q},
\]
using \code{np.linalg.det} for $U$ and \code{np.linalg.solve} to form
$\Lambda^{-1}N$ without explicitly inverting $\Lambda$. This is exactly the
$Q_{\text{eff}}$ that, multiplied by $D$, becomes the $\mathcal{B}$ fed to the
production transform of \S\ref{sec:hk-loops} --- the oracle and the production
path compute the \emph{same} object two different ways.

\subsection{\code{temporal\_integrate.py}: the causal time integral}

After the momentum integral, the residual integral is over the edges' time
(Schwinger) parameters with the causal structure. For a 2-vertex self-energy
where all internal edges span the same inter-vertex time $t$,
\code{sigma\_parametric} (\file{temporal\_integrate.py:56}) assembles
\[
  \Sigma(q,t) \;=\; T^{n_C}\int_{[|t|,\infty)^{n_C}}\!\prod\dd s_C\;
  \ee^{-\mu\sum_e w_e}\;I_{\text{mom}}\big(\{(a_e,b_e)\},\{w_e\},q\big),
\]
where a \emph{retarded} edge has fixed duration $w_e=t$ (from $\theta(t)$), and a
\emph{correlation} edge carries a Schwinger parameter $w_e=s_e$ integrated over
$[|t|,\infty)$, each contributing a factor $T$ (temperature). The combinatorial
factor $\Sym(\Gamma)$ is applied by the caller, not here. The dispatch on the
number of correlation edges $n_C$ is the interesting part: $n_C=0$ evaluates
directly, $n_C=1$ uses an adaptive \code{scipy.integrate.quad}, and $n_C\ge 2$
uses a \emph{Gauss--Laguerre tensor rule} over the correlation parameters.

\begin{note}
\textbf{Gauss--Laguerre quadrature} integrates $\int_0^{\infty}\ee^{-x}g(x)\,\dd x$
\emph{exactly} for polynomial $g$ up to degree $2\,\text{deg}-1$, using the nodes
and weights returned by \code{np.polynomial.laguerre.laggauss(deg)}. Because the
Schwinger integrand carries an $\ee^{-\mu s}$ weight, the substitution
$s_C = lo + x_k/\mu$ turns each correlation-edge integral into exactly that form,
and the $n_C\ge 2$ case takes the Cartesian product of node indices with
\code{itertools.product}. This is roughly 10--100$\times$ faster than the
adaptive \code{dblquad}/\code{nquad} it replaces --- the two-correlation-edge
\code{dblquad} in \code{loop\_parametric.sigma\_K\_kernel} was the bottleneck of
the older C stack.
\end{note}

The convenience edge-specs \code{bubble\_edges} and \code{sunset\_edges}
(\file{temporal\_integrate.py:127,135}) supply the routing triples for the
validated 2-vertex topologies (the 1-loop bubble and the 2-loop sunset), and
\code{bubble\_delta\_equal\_time\_via\_C} (\file{temporal\_integrate.py:144})
collapses the tabulated self-energies into $\delta C(q,0)$ via the \msrjd{}
Dyson equation. At $d=1$ this reproduces the backend-B golden reference; at
$d=2,3$ it gives the higher-dimensional bubble. Again: this is the
\emph{cross-check}, kept to certify that the production integrator produces the
same numbers.

\section{Guards, errors, and known limits}\label{sec:hk-guards}

The subsystem is built to \emph{refuse} cases it cannot handle correctly rather
than return a plausible-looking wrong number. Every refusal is a
\code{SpatialPropagatorError} (\file{heat\_kernel.py:50}) --- the subsystem's one
exception type, raised whenever the Tier-1 closed-form path does not apply ---
or, downstream, a clear \code{NotImplementedError}. Collected in one place, the
guards are:

\begin{description}
  \item[The diagonal (Tier-1) restriction.] The closed-form heat kernel requires
        a \emph{diagonal} inverse propagator (each field an independent
        Allen--Cahn mode). Any off-diagonal coupling makes
        \code{build\_spatial\_propagator} raise
        (\file{heat\_kernel.py:369--377}) and leaves \code{G\_tx\_sym = None};
        \code{compute\_spatial\_correlator\_tree} then raises a clear
        \code{NotImplementedError} (\file{spatial\_correlator.py:239--248}), and
        the bridge re-routes a coupled model to the spectral driver of
        Chapter~\ref{ch:spatial-coupled}.
  \item[Higher-derivative operators are rejected.]
        \code{extract\_mass\_diffusion} rejects any $k$-power $>2$
        (\file{heat\_kernel.py:213--216}) --- a $(\Lap)^{2}$ (a $k^{4}$, e.g.\ a
        conserved Model-B \emph{kinetic} operator) is a v2 feature. (Conserved,
        $q^{2}$-dependent \emph{noise} is separately skipped by the noise
        extractors, left at zero.)
  \item[The $\ii\omega$ normalization is strict.] The time-derivative coefficient
        must be exactly $\ii$ (\file{heat\_kernel.py:201}). A model that hides a
        factor on $\Dt$ is rejected as ``not a standard \msrjd{} kernel'' rather
        than silently rescaled.
  \item[Diffusion sign and zero.] \code{compute\_spatial\_correlator\_tree}
        raises on $B<0$ (anti-diffusive, ill-posed, the heat kernel diverges ---
        check the sign of $-D\,\Lap$ with $D>0$) and on $B=0$ (no Laplacian: a
        time-only field, for which a spatial correlator is undefined)
        (\file{spatial\_correlator.py:261--275}). The underlying reason is that
        the erf form divides by $\sqrt{4\pi B}$ and takes $\sqrt{}$ of the mass.
  \item[Periodic boundaries and drift are $d=1$ only.] \code{image\_sum} raises
        for $d\ne 1$ (\file{heat\_kernel.py:143--145});
        \code{gaussian\_heat\_kernel} raises for $V\ne 0$ with $d\ne 1$
        (\file{heat\_kernel.py:85--88}); and \code{build\_spatial\_propagator}
        raises for periodic boundaries with $d\ne 1$
        (\file{heat\_kernel.py:345--349}).
\end{description}

The v2 backlog these guards delimit is, in one list: off-diagonal (coupled
multi-field) propagators, higher-derivative operators ($k^{4}$+ kinetic terms),
and $d\ge 2$ derivative vertices (the one remaining numerical-transform gap of
\S\ref{sec:hk-loops}).

\section{What this chapter covered}

The spine of this chapter is one identity --- \emph{the inverse transform of a
Gaussian is a Gaussian} --- applied twice. Closing the $\omega$-contour on the
diffusive inverse propagator $A + Bk^{2} + \ii\omega$ gives exponential decay in
time, $\theta(t)\,\ee^{-(A+Bk^{2})t}$; inverse-transforming the surviving
momentum Gaussian gives the real-space heat kernel
$\theta(t)\,(4\pi Bt)^{-d/2}\,\ee^{-x^{2}/4Bt - At}$, the function
\code{gaussian\_heat\_kernel}. From there, drift shifts the Gaussian centre to
$x=vt$ (and is usually zero at the homogeneous saddle); a periodic box is the
super-exponentially convergent image sum \code{image\_sum}; the three coefficients
$(A,B,V)$ are read symbolically off \code{K\_ft} with Sage by
\code{extract\_mass\_diffusion}; the picklable propagator block is assembled by
\code{build\_spatial\_propagator} and its non-picklable callables rebuilt by
\code{make\_g\_tx\_callables}; the tree-level correlator is noise driving two
retarded legs, collapsed to the mpmath-precision erf integral
\code{erf\_time\_integral} inside \code{free\_two\_point}; and the noise strength
itself is read out of the action's $(2,0)$ sector by
\code{extract\_noise\_coefficients}.

At loop order the same Gaussian transform runs batched and vectorized in
\code{full\_integrator.\_heat\_kernel\_x\_general}, gated by \code{\_use\_analytic}
(plain vertices at any $d$, or any vertices at $d=1$), with a numerical-transform
cross-check (\code{radial\_inverse\_ft}, \code{periodic\_inverse\_ft}) kept alive
behind \code{SPATIAL\_FORCE\_NUMERICAL\_FT} and friends, and pinned by the
oracle-only Symanzik modules \code{loop\_parametric.py} and
\code{temporal\_integrate.py}. The next chapter,
Chapter~\ref{ch:spatial-coupled}, takes up the case this one explicitly refuses
--- the coupled, off-diagonal multi-field propagator --- where a single diagonal
heat kernel no longer suffices and the spectral construction takes over.
